\documentclass[11pt]{article}

\usepackage[a4paper,margin=0.85in]{geometry}
\usepackage{setspace}
\usepackage{graphicx}
\usepackage{array}
\usepackage{parskip}
\usepackage{ragged2e}

\pagestyle{empty}
\setstretch{1.05}

\begin{document}

\begin{flushright}
Zhongli Wang \\
+49 160 9541 6015 \\
nkwangzhongli@gmail.com \\
Astastrasse 22, 33617 Bielefeld\\ Germany \\
\end{flushright}

\vspace{0.1cm}

\noindent
Recruiting Team \\
Center for Securities Analysis with Financial Technology \\
Department of Accounting \\
The Hong Kong University of Science and Technology \\
Hong Kong \\

\vspace{0.1cm}

\noindent
\textbf{Application for Post-doctoral Fellow} \\

\vspace{0.05cm}

\noindent Dear Recruiting Team,

I am currently completing my Ph.D.\ in Economics at Bielefeld University and Universit\`a Cattolica del Sacro Cuore, with planned completion in September 2026. My doctoral research combines empirical analysis, machine learning, statistical modeling, agent-based simulation, and financial market analysis. I am very interested in the Post-doctoral Fellow position at the Center for Securities Analysis with Financial Technology because it closely matches my background in financial data analysis, fintech-related methods, machine learning, and research on firm behavior, financial irregularities, and market dynamics.

My doctoral work develops agent-based models of market environments where outcomes are non-linear, path-dependent, and shaped by feedback loops. In my dissertation, \emph{Strategic Behavior and Economic Dynamics in ICT Markets: An Agent-Based Framework}, I study how firms and heterogeneous users interact under network effects, congestion constraints, imperfect information, and strategic investment decisions. Methodologically, this work required me to formalize behavioral and strategic decision rules, build scalable simulation pipelines, conduct counterfactual analysis, and evaluate how different market or policy conditions affect firm outcomes. These skills are highly relevant to securities analysis and financial technology research, where firm performance, investor behavior, information frictions, market structure, and regulatory or technological changes often interact in complex and adaptive ways.

During my doctoral studies, I have developed end-to-end analytical workflows for complex economic and financial questions. This includes data collection, cleaning, validation, database construction, feature engineering, exploratory analysis, model estimation, robustness checks, and the preparation of structured research outputs. In one dissertation project, I estimate an agent-based model of competition and investment dynamics using financial and operational data from AT\&T and T-Mobile, applying simulation-based estimation and Python-based analytical pipelines. This project strengthened my ability to work carefully with firm-level data, evaluate assumptions, maintain reproducible research workflows, and translate quantitative outputs into economic interpretation.

A second line of my work is directly connected to financial-market analysis and fintech-related research. In collaborative research on \emph{Polytope Fraud Theory}, I use supervised machine-learning methods to detect structured anomalies in financial statements, with a focus on firms targeted by activist short sellers. This work involved constructing firm-level features, handling heterogeneous financial indicators, evaluating model performance through recall, F1 score, predictive accuracy, and feature importance, and interpreting algorithmic signals in terms of financial misconduct, investor protection, and market transparency. This experience is particularly relevant to the Center's work on securities analysis, as it connects machine learning, financial statement information, anomaly detection, and economically meaningful interpretation of market-relevant signals.

I can contribute to the Center's research activities by connecting three components into a coherent empirical research pipeline: first, the construction and maintenance of clean financial and securities-related datasets; second, the application of statistical and machine-learning methods to identify patterns, anomalies, and predictive signals; and third, the translation of quantitative results into research outputs suitable for academic papers, working papers, and policy- or practitioner-oriented communication. I am comfortable working with messy and heterogeneous data sources, documenting extraction and coding procedures, checking data consistency, and building reproducible workflows so that results remain auditable and reusable.

My technical background is well aligned with the requirements of the position. I have strong experience in Python, including pandas, NumPy, scikit-learn, XGBoost, statistical analysis, machine learning, simulation, and reproducible data-processing pipelines. I also have experience with Stata for empirical analysis and am comfortable adapting to additional statistical software such as SAS when required. In addition, I use SQL for analytical querying, joins, aggregation, and data validation, and I am continuing to strengthen my database-management skills for large-scale financial and accounting datasets. I have also worked in Linux-based research environments and used structured project workflows for simulation and data analysis.

I am also interested in the growing role of large language models and agentic workflows in empirical research. In a related project, \emph{Agentic Workflows for Economic Research}, I helped develop a framework for integrating large-language-model agents into research pipelines, including document collection, classification, workflow documentation, and research assistance. For research in securities analysis and financial technology, such tools can be useful for collecting and organizing textual information, supporting literature review, processing firm disclosures, and improving transparency in data construction and documentation. I would be excited to further develop these skills in support of the Center's research agenda.

Beyond technical skills, I bring strong intellectual curiosity, responsibility, and experience in academic communication. My Ph.D.\ training has required me to write research papers, prepare presentations, document methodological choices, and communicate complex quantitative findings clearly in English. I value timely communication, careful documentation, and collaborative research practice. I would be glad to assist faculty members with data collection and cleaning, empirical analysis, database management, literature and editorial support, research documentation, and the preparation of materials for academic publication.

I am particularly attracted to HKUST because of its strong research environment and its location in Hong Kong, where financial markets, technology, and international academic exchange are closely connected. The Center for Securities Analysis with Financial Technology offers an excellent environment for applying my background in economics, machine learning, and financial data analysis to research questions with both academic and practical relevance. I would be highly motivated to contribute to the Center's projects while deepening my expertise in accounting, securities analysis, and financial technology.

Thank you so much for considering my application! I would be appreciated to discuss how my background and abilities can contribute to the research activities of the Center for Securities Analysis with Financial Technology at HKUST.

\vspace{0.15cm}

\noindent Sincerely, \\[0.8cm]
\noindent Zhongli Wang

\end{document}
